% JAC1964-20261121. /root/programme_audit; GPT-6 Astra Pro.
% Overview only; root owns integration and the master proof text.

\section{The seven-theorem programme: interpretations and scope}
\label{sec:programme-overview}

The operative programme in cote~44 is on scans~2, 4 and~6.
The letter to Serre dated 12~August~1964 occupies scans~8--9.
Scan~1 is a modern cover; scans~3, 5 and~7 concern graded
completion or root-system material and are not further claims in
this programme. The supplied sources are TeX transcriptions;
the manuscript images were unavailable. The following is a
mathematical audit, not a palaeographic verification. Uncertain
readings and deleted clauses remain part of the source record.

Throughout the table, $P^p(X)$ means Saito's higher Picard variety
of algebraically trivial cycles modulo incidence equivalence.
This is the explicit interpretation under which the programme is
tested. A universal regular algebraic representative is identified
with it only in the cases where the comparison has been proved
separately. Polarizations and duality statements are rational;
no integral universal cycle or principal polarization is asserted.

\begingroup
\small
\setlength{\tabcolsep}{4pt}
\renewcommand{\arraystretch}{1.15}
\begin{longtable}{@{}p{0.08\linewidth}p{0.19\linewidth}p{0.68\linewidth}@{}}
\toprule
Claim & Exact archive locus & Interpretation and result\\
\midrule
\endfirsthead
\toprule
Claim & Exact archive locus & Interpretation and result\\
\midrule
\endhead
Th.~1 & Cote~44, scan~2; clarified by scan~8 &
The family map from $\operatorname{Alb}(S)$ has the incidence
kernel, as a statement on geometric points. Saito's construction
supplies this reading; the uncertain archive notation is not
silently replaced by universal regular equivalence.\\

Th.~2 & Cote~44, scan~2; scan~8, Betti bound &
$\dim P^p(X)\le b_{2p-1}(X)/2$. This is established background for
smooth projective varieties over algebraically closed fields in
every characteristic: Saito, Proposition~3.4, pp.~109--110.\\

Th.~3 & Cote~44, scan~2; scan~8, complementary indices &
The incidence quotient is an abelian variety and $P^p(X)$ is
canonically dual, up to isogeny, to $P^{n+1-p}(X)$.
The index $n+1-p$ is forced by the pairing and is explicit in
scan~8. Saito, Theorem~3.6 and Proposition~4.4, supplies the
all-characteristic construction and duality.\\

Th.~4 & Cote~44, scan~2; scan~8, primitive margin; scan~9, formula for $D$ &
The correct signed class on the primitive middle factor is
$(-1)^m\Delta^*d/2$ for $\dim X=2m-1$.
JAC-T4 proves ampleness for abelian $X$ in every characteristic
under 032, Corollary~8.4. The generic-section and complete-intersection
extensions are H9-T6a and H9-T13. The literal unsigned assertion
on the full factor fails already on $E^3$ (JAC-P03).\\

Th.~5 & Cote~44, scan~4, height formula and fibre hypothesis &
The source assumes algebraically trivial restrictions to the generic
section. Its uncorrected formula fails even for geometrically
integral quadric-pencil fibres (H9-C1). H9-T1 gives the corrected
intersection formula; H9-T9 identifies the signed height with a
positive form modulo trace on a polarized algebraically trivial
primitive incidence domain. The larger homological domain requires
the qualifications below; whole-group positivity after a separable
cover is false (H9-T17).\\

Th.~6 & Cote~44, scan~4, three degree ranges &
For $\dim X=n$, restriction is an isogeny when $2p\le n-1$,
its image is the fixed factor when $2p=n$, and it is surjective
when $2p\ge n+1$. JAC-P02 proves these rational assertions when
the normalized Lefschetz lowering operator of $X$ is algebraic,
in particular for abelian $X$ in every characteristic. JAC-P16
extends this to every variety of geometric abelian type, using
032's numerical/homological equality and Jannsen semisimplicity. It also
proves the finite-kernel assertion deleted in the critical clause;
the deletion is not erased from the source.\\

Th.~7 & Cote~44, scan~6 &
With customary $\tau$-equivalence, a nonzero multiple is
algebraically equivalent to zero. The assertion is false in
characteristic zero even for the CM Fermat-quartic Jacobian
(JAC-P05, using Bloch's non-torsion theorem).
It is true for every abelian variety over $\overline{\mathbf F}_p$,
indeed every geometric abelian-type variety, by JAC-P15 using
032, Corollary~8.3, and Kahn's theorem. H9-T14 further proves
numerical=$\tau$ after any extension of the constant field for
these constant varieties, and for abelian varieties geometrically
isogenous to such constants. Divisors in every
characteristic and zero-cycles on connected smooth projective
varieties are classical positive cases.\\
\bottomrule
\end{longtable}
\endgroup

\paragraph{The height domains are part of the assertion.}
Replacing the algebraic-triviality hypothesis of Th.~5 by
homological triviality changes the problem. JAC-P04 exhibits that
distinction on an abelian fourfold. For a genuine Lefschetz pencil
on an even-dimensional abelian variety, JAC-P07 proves condition~A
automatically in every characteristic. H9-T7 then computes the
height on the entire homologically trivial Chow group as an
ambient primitive intersection form plus a variable-axis form,
and identifies the remaining radical with the radical of the axis
form. This is positive in the surface-axis case
H9-T8 and for the geometric generic ordered axes of degree
$d\ge3$ in JAC-P09, using 032 for the ambient term. The latter
scope is the original pencil field with a geometric generic axis.
In characteristic zero, classical Hodge--Riemann positivity proves
the sign for every smooth axis and every $m$ over the original
pencil field (H9-T8a), without 032. Thus the remaining arbitrary-axis
sign question for that field is in positive characteristic.
H9-T12 proves the precise finite-extension statements: the whole
result persists under purely inseparable extensions; arbitrary
finite extensions have an orthogonal restriction--norm
decomposition on the admissible domain; the algebraically trivial
primitive domain remains positive modulo the unchanged fixed trace.
Whole-group semipositivity after a separable cover is refuted by
H9-T17 in every characteristic. On an $E^4$ pencil after a
separable elliptic cover unramified above the critical values,
the admissible norm-zero class $hD$ has exact height $-12d^3$.
One may choose the cover degree $N\ge3$ prime to the
characteristic; fibre divisors linearly equivalent to $N\cdot0$
give the exact norm identity. In characteristic zero and generic
dimension three, Zhang's theorem instead makes the
Albanese-primitive admissible part positive semidefinite; the
orthogonal Lefschetz part $h\operatorname{Pic}^0(Y)(K)_{\mathbf Q}$
is negative semidefinite with exact trace kernel.
For constant abelian varieties descending to
$\overline{\mathbf F}_p$, H9-T14 identifies the entire homologically
trivial domain with the algebraically trivial domain, even after
transcendental extension. Its primitive height quotient is
$\operatorname{Hom}_k(J(C),B)_{\mathbf Q}$, with radical the
inverse image of $B(k)_{\mathbf Q}$. The isotrivial version uses
the actual $K/k$-trace instead of assuming that the variety is
already constant.

\paragraph{The primitive and representative qualifications are essential.}
The margin in cote~44, scan~8, explicitly requires the primitive
decomposition. Under $\phi_L(x)=t_x^*L\otimes L^{-1}$, the
elliptic normalized diagonal induces the negative of the ample
polarization. Thus the divisor-incidence sign is $(-1)^m$,
whereas the primitive odd cup-product sign is $(-1)^{m-1}$.
The factor $1/2$ is the diagonal extraction of a normalized
symmetric divisor correspondence. H9-T10 proves the comparison
with the actual regular representative $\operatorname{Ab}^2$
for the stated generic threefold sections; H9-T13 extends that
case to the indicated complete intersections in an abelian variety.
Neither result supplies a universal regular representative in
arbitrary codimension.

\paragraph{The finite-field conclusion uses the stronger accepted input.}
The Hodge-standard assertion of 032, Corollary~8.4, is not by
itself a rational-equivalence theorem. JAC-P15 uses its Tate
consequence, Corollary~8.3, printed pp.~49--50, together with
Kahn, Theorem~1.10, printed p.~982. It proves
$\operatorname{CH}^r_{\mathbf Q}=N^r_{\mathbf Q}$ over finite
fields and their algebraic closures for abelian-type varieties.
This is compatible with the characteristic-zero Ceresa
counterexample and does not assert rational=numerical over
arbitrary transcendental constant-field extensions.
H9-T14 instead proves numerical=homological=algebraic equivalence
with rational coefficients on every such constant-field
extension. Its direct spreading argument compares intersection
degrees with all complementary constant cycles and therefore
requires no additional standard conjecture on the extended
variety. This is the customary numerical--$\tau$ conclusion;
it does not assert that the algebraically trivial Chow group vanishes.

\subsection*{Finding-to-archive atlas}

The IDs below name deductions or calculations made in this run.
The archive column identifies the claim they resolve or the
passage that motivates them; it does not attribute their modern
inputs or their proofs to Grothendieck. Saito's construction,
the SGA monodromy theorems, Bloch's non-torsion theorem, Kimura
nilpotence, and Kahn's Tate--Beilinson theorem remain cited
background.

\begingroup
\small
\setlength{\tabcolsep}{4pt}
\renewcommand{\arraystretch}{1.12}
\begin{longtable}{@{}p{0.12\linewidth}p{0.25\linewidth}p{0.58\linewidth}@{}}
\toprule
ID & Grothendieck locator & Finding and precise role\\
\midrule
\endfirsthead
\toprule
ID & Grothendieck locator & Finding and precise role\\
\midrule
\endhead
JAC-P00 & Cote~44, scans~2 and~8; Th.~6 on scan~4 &
Rational descent of the incidence quotient and its normalized
families, including imperfect pencil fields and inseparable
norm multiplicities.\\
JAC-P01 & Cote~44, scan~2, Ths.~1--3; scan~8, Tate-module comparison &
The relevant odd realization detects incidence homomorphisms;
cohomological algebraic projectors induce actual rational
abelian-variety projectors.\\
JAC-P02 & Cote~44, scan~4, Th.~6; cote~16, scans~40 and~49 &
Restriction, critical fixed trace, and surjectivity under algebraic
normalized lowering. The critical trace proof spreads an actual
cycle and uses the blowup formula.\\
JAC-P03 & Cote~44, scan~8, primitive margin; scan~9, $D$ &
Explicit families on $E^3$ give opposite signs. The calculation
fixes the divisor-incidence sign and shows that quotienting the
kernel alone does not produce positivity.\\
JAC-P04 & Cote~44, scan~4, algebraic-triviality hypothesis &
An explicit primitive cycle on $E^4$ has homologically trivial
but non-algebraically-trivial restriction in the stated high-degree
characteristic-zero setting.\\
JAC-P05 & Cote~44, scan~6, Th.~7 &
Bloch's infinite-order Fermat Ceresa class refutes customary
numerical=$\tau$ even for a CM abelian threefold in characteristic zero.\\
JAC-P06 & Cote~44, scan~6, undefined $\tau$ &
Over $\overline{\mathbf F}_p$, algebraically trivial Chow classes
are torsion, so the numerical--$\tau$ claim is a rational-Chow
comparison. JAC-P15 supplies that comparison in abelian type.\\
JAC-P07 & Cote~44, scan~4, pencil/fibre hypothesis in Th.~5 &
For an abelian $n$-fold, $n$ even, variable rank is
$h^n-\binom{2n}{n}/(n+1)>0$. This proves condition~A and
the required specialization injectivity for genuine Lefschetz
pencils, in every characteristic.\\
JAC-P08 & Cote~44, scan~4, Th.~6; cote~16, scans~40 and~49--50 &
Explicit algebraic lowering for smooth complete intersections,
and stability under products, extend the proved restriction/trace
statements beyond abelian ambient varieties.\\
JAC-P09 & Cote~44, scan~4, Th.~5; enlarged H9 domain &
Orthogonal monodromy and finite orbit of a geometric generic
cycle eliminate the variable algebraic axis classes for ordered
degree-$d$ axes, $d\ge3$, in every characteristic.\\
JAC-P10 & Arithmetic continuation of cote~44, scan~4, height question &
The explicit stable Hodge lattice of the Fermat quartic has
normalized determinant gain $7/4$; hence
$h_F(J)-3h_F(E_{1728})=\log(2)/4$. No node thickness is inferred
from an unweighted intersection matrix.\\
JAC-P11 & Cote~44, scan~6, non-$\tau$ cycles; arithmetic height continuation &
The central CM operation $T$ has $T^2=-1$ and preserves height.
The Ceresa vector and $T$-translate span a positive orthogonal
plane and are independent even modulo algebraic equivalence.\\
JAC-P12 & Same Th.~7 scope question &
For the hyperflex embedding, the ordinary tautological
codimension-two ring is $\mathbf Q\theta^2+\mathbf Q\eta$;
$T\eta$ lies outside it even modulo algebraic equivalence.
Tautological generation does not prove generation of the whole Chow group.\\
JAC-P13 & Cote~44, scan~6; modern motivic continuation &
With explicitly added Kimura background, the exact complement is
$N=h^3(J)^G$. Its whole $\operatorname{CH}^2$ is a
$\mathbf Q(i)$-module carrying the Hermitian height form.\\
JAC-P14 & Same modern continuation &
For $\operatorname{End}^0(J)=M_3(\mathbf Q(i))$, the exact action
on $N$ is $u-vT$ when $\det b=u+vi$. The full abelian-endomorphism
orbit of the Ceresa vector remains its CM plane.\\
JAC-P15 & Cote~44, scan~6, Th.~7 over finite fields &
Accepted full Tate plus Kahn gives rational=numerical, and hence
numerical=$\tau$, for every geometric abelian-type variety over
finite fields or $\overline{\mathbf F}_p$.\\
JAC-P16 & Cote~44, scan~4, Th.~6; cote~16, scan~49 &
Numerical=homological on all powers of a geometric abelian-type
variety, with algebraic Kunneth projectors and hard-Lefschetz
inverses over the original field. This proves $B(X)$ and extends
Th.~6 in every characteristic; it does not assert positivity of
$X$'s own primitive form.\\
JAC-P17 & Arithmetic continuation of cote~44, scan~4 &
The exact Fermat motive and Ceresa vector descend to $\mathbf Q$;
the central operation $T$ is defined over $\mathbf Q(i)$ and is
anti-invariant under conjugation. The entire rational Chow group
over $\mathbf Q(i)$ is the orthogonal sum $V_{\mathbf Q}\oplus
T V_{\mathbf Q}$, with no generation assertion.\\
JAC-P18 & Cote~44, scan~6; arithmetic radical continuation &
A nonzero algebraically trivial class on the Fermat summand
produces a nonzero homologically trivial cycle in the degree-four
component of $\operatorname{CH}^2(C\times J)$ for some curve $C$
after a finite field extension. Murre~(D) on this component would
eliminate that algebraic radical; no converse or full
Abel--Jacobi injectivity is asserted.\\
\bottomrule
\end{longtable}
\endgroup

\paragraph{The relevant cote~16 corrections are mathematical, not palaeographic.}
In no.~7, scans~34, 36, 38, 40, 49 and~50 are the operative
Lefschetz pages; scans~35, 37 and~39 are unrelated printed
material. The lowering operator is normalized along a Lefschetz
string, not the operator with $\mathfrak{sl}_2$ string
coefficients. With that convention the boxed identity on scan~40,
\[
             \Lambda_Y=j^*\Lambda_X^2j_*,
\]
is correct and is used in JAC-P02. Several adjacent indices in
the transcription are type-inconsistent: for example the lower
right target in scan~34 is forced by degree to be
$H^{2n-i}(Y)$, and the first intermediate identity on scan~40
has only the bottom-endpoint sum $0\le i\le n$ in
$L_X\Lambda_X=1-\sum p_i^X$. The full all-endpoint sum there
belongs only to the subsequent product. The detailed
transcription audit retains every uncertainty and distinguishes
these forced mathematical readings from claims about the
unseen manuscript.
